% ECONOMICS_PROBLEM: {"id": "O32-593", "title": "Scientific discovery bottlenecks", "jel_code": "O32", "source_id": "OP-0410"}
% !TeX program = xelatex
\documentclass[11pt,a4paper]{article}
\usepackage[margin=23mm]{geometry}
\usepackage{fontspec}
\setmainfont{Latin Modern Roman}
\usepackage{amsmath,amssymb,mathtools}
\usepackage{xcolor,enumitem,booktabs,longtable,array,xurl,hyperref}
\definecolor{muted}{HTML}{53636C}
\hypersetup{colorlinks=true,urlcolor=blue,linkcolor=blue}
\setlength{\parindent}{0pt}
\setlength{\parskip}{5pt}
\newcommand{\R}{\mathbb R}
\newcommand{\E}{\mathbb E}
\newcommand{\N}{\mathbb N}
\newcommand{\ind}{\mathbf 1}
\newcommand{\Var}{\operatorname{Var}}
\newcommand{\Cov}{\operatorname{Cov}}
\newcommand{\supp}{\operatorname{supp}}
\DeclareMathOperator*{\argmin}{arg\,min}
\DeclareMathOperator*{\argmax}{arg\,max}
\newcommand{\entrymeta}[3]{{\sffamily\small\color{muted}\textbf{\texttt{#1}}\quad|\quad #2\par #3\par}}
\newcommand{\Problemref}[1]{\texttt{#1}}
\newcommand{\JELref}[2]{\texttt{#2} (JEL \texttt{#1})}
\begin{document}
\section*{Scientific discovery bottlenecks}
\textbf{Problem ID:} \texttt{O32-593}\quad \textbf{JEL:} \texttt{O32}\par
% BEGIN ECONOMICS STATEMENT
\label{problem:OP-0410}
{\sffamily\small\color{muted}Original subject group: Innovation and AI economics\par}
\label{definitions:problem:30007692}
\textbf{Audit classification:} Explicit published open question. The source explicitly poses discovery-bottleneck questions, but the pharmaceutical queue experiment is editorial. Discrete capacity changes need not have derivatives.
\subsection*{Definitions and precise research target}
Consider pharmaceutical discovery projects entering a prespecified laboratory network during one year. Each project passes through hypothesis generation \(H\), experimental testing \(E\), and independent validation \(V\), with staff hours \(c_s\) and processing capacity \(k_s\) at stage \(s\in\{H,E,V\}\). An intervention \(a=(a_H,a_E,a_V)\in\{0,1\}^3\) supplies a fixed AI assistance package at selected stages; zero denotes the existing workflow. Let \(N(a)\) be the number of discoveries meeting a blinded, preregistered reproducibility and clinical relevance threshold within twelve months, and let \(C(a)>0\) include labor, computation, experiments and validation costs. For the distribution of eligible projects and laboratories, define productivity \(P(a)=E[N(a)]/E[C(a)]\). A stage is a bottleneck for a specified feasible capacity increment if it gives the largest increment of expected validated discoveries per additional dollar; use derivatives only when they exist. Estimate the factorial intervention effects on \(P(a)\), stage waiting times and marginal capacity values, allowing faster hypothesis generation to congest testing or validation. Laboratory randomization, recorded queues and independent outcome adjudication are proposed inputs; within-laboratory interactions must be retained. Which stages remain binding, and when does investment in a downstream stage create larger validated gains than further AI improvement upstream? The target is this laboratory population, rather than all scientific discovery.

\textbf{Definitions, assumptions and mathematical qualifications.} Treat each laboratory as a cluster containing a fixed intake cohort of projects. \(N(a)\) counts eligible projects reaching a preregistered validation threshold within twelve months; failed, delayed and abandoned projects count zero. \(C(a)\) is total real resource expenditure for that cohort, with \(0<c_{\min}\leq E[C(a)]<\infty\) and finite \(E[N(a)]\). The ratio of expectations differs from \(E[N(a)/C(a)]\). Let \(n(k,a)\) be expected completed discoveries for capacity vector \(k\), and \(g_s(\Delta)>0\) the cost of a feasible capacity increment. Define finite-increment bottleneck value \(b_s(\Delta)=\{n(k+\Delta e_s,a)-n(k,a)\}/g_s(\Delta)\). Derivatives are used only when limits exist; queues, prerequisites and project allocation are part of the experimental treatment.
\subsection*{Variants and assumption changes}
\begin{enumerate}
\item \textbf{Editorial, factorial deployment.} Randomize all eight AI stage packages across laboratories and identify productivity, waiting-time and reproducibility contrasts for the specified cohort.
\item \textbf{Editorial, capacity optimization.} Also randomize feasible testing or validation capacity increases; compare \(b_s(\Delta)\) under a fixed resource budget, accounting for congestion created by upstream acceleration.
\end{enumerate}
\subsection*{References and source audit}
\textbf{Checked source.} 2024 Stanford draft, Section 3.2, PDF pp. 10--11. Published questions concern discovery constraints, disciplinary boundaries and access. Full primary draft and cover retrieved. \url{https://digitaleconomy.stanford.edu/app/uploads/2024/11/ETAI-White-Paper.pdf}.
\begin{enumerate}\raggedright
\item\hypertarget{definitions:source:30007692:1}{} Erik Brynjolfsson; Anton Korinek; Ajay Agrawal. 2024. \emph{The Economics of Transformative AI: A Research Agenda}. 2024 Stanford draft, Section 3.2, PDF pp. 10--11.
\url{https://digitaleconomy.stanford.edu/app/uploads/2024/11/ETAI-White-Paper.pdf}.
\end{enumerate}

\subsection*{Common conventions: Source mathematical conventions}

These conventions apply to each entry unless explicitly replaced.

\textbf{Models and identification.} A primitive model \(m\) specifies all probability laws, feasible choices, information, equilibrium and measurement rules used by its target. The admissible class \(\mathcal M\) is fixed independently of outcomes. If \(P_O(m)\) is its observable probability law and \(\tau(m)\) the target, its identified set at observable law \(P\) is
\[
 \mathcal I_\tau(P)=\{\tau(m):m\in\mathcal M,\ P_O(m)=P\}.
\]
Point identification means that this set is a singleton. An empty set signals incompatibility of data and assumptions. Coordinate bounds need not describe the joint set. Bounds are sharp only if every retained target is attainable or arbitrarily approachable by compatible models. Finite bounds require explicit restrictions on unobserved quantities; unrestricted confounding, hidden wealth, extrapolation or arbitrary equilibrium selection can yield uninformative sets.

\textbf{Causal interventions.} A potential outcome \(Y_i(p)\) is the outcome under a complete policy regime \(p\), including timing, financing, information and market spillovers. The target population and units are fixed before treatment. With interference, use the complete assignment vector or a justified exposure mapping; individual no-interference notation alone cannot identify an economy-wide intervention. A difference of mean potential outcomes does not require the cross-world joint distribution. Conditioning on post-treatment migration, survival, enrollment or employment changes the estimand and needs separate justification.

\textbf{Statistical statements.} An estimator, confidence set, forecast or test specifies a sampling law, model class and loss/coverage criterion. Uniform \((1-\alpha)\)-coverage means \(\inf_{m\in\mathcal M}\Pr_m\{\tau(m)\in C_n\}\geq1-\alpha\), or an explicitly asymptotic version. The whole parameter space trivially has coverage, so informativeness requires a stated width, risk or power objective. A forecasting improvement is relative to a named benchmark, information set and evaluation population.

\textbf{Equilibrium and optimization.} An equilibrium comprises feasible plans, optimality, beliefs where needed, budget constraints and all market/resource clearing. The primitives or a stated benchmark define each correspondence \(\mathcal E(p;m)\). Nonemptiness is assumed or proved, never inferred just from compact policy sets. A selected-equilibrium target uses a declared selection rule; a robust target quantifies over all permitted equilibria. Use suprema or infima unless attainment is established. An \(\varepsilon\)-optimal policy is feasible and within \(\varepsilon\) of the finite optimum value. Report an empty feasible set separately.

\textbf{Welfare and accounting.} Normative utility, interpersonal normalization, social weights, numeraire, discounting and the use of public receipts are primitives. Behavioral utility need not coincide with the welfare criterion. Household welfare includes ownership income and rents once; adding profits again to compensating variation generally double-counts them. A monetary equivalent is defined by an explicit transfer and price convention, with monotonicity and range conditions. Average effects, mechanism contributions, transfers and welfare losses are different targets. Infinite sums require convergence and dynamic feasible plans require terminal or no-Ponzi conditions.

\textbf{Source versions.} A working-paper date, revision date and journal publication year are distinguished. A DOI for a journal version is not automatically the DOI of an earlier manuscript. Locators refer to the stated version and printed pages where specified. A metadata-only check does not verify a claim about a full-text conclusion. Every entry's source subsection narrows the provenance claim accordingly.
% END ECONOMICS STATEMENT
\end{document}
